\chapter{Introduction}
\label{ch:into}

This chapter presents the project context and objectives. It describes the problem
being investigated, outlines the aims and objectives, and summarises the solution
approach and project management plan.




%%%%%%%%%%%%%%%%%%%%%%%%%%%%%%%%%%%%%%%%%%%%%%%%%%%%%%%%%%%%%%%%%%%%%%%%%%%%%%%%%%%
\section{Project Presentation}
\label{sec:proj_pres}

\subsection{Background}
\label{sec:into_back}

Reinforcement Learning (RL) has emerged as one of the most promising paradigms 
in artificial intelligence, enabling autonomous systems to learn complex 
behaviours from raw experience without explicit programming. Landmark achievements 
such as DeepMind's AlphaGo~\citep{silver2016mastering}, and the DRL agent 
playing Atari games at superhuman level~\citep{mnih2015humanlevel}, have 
demonstrated that RL agents can master tasks previously thought to require 
human intuition.

Multi-agent reinforcement learning (MARL) extends these ideas to environments 
where multiple agents learn simultaneously, introducing additional challenges 
such as non-stationarity, coordination, and emergent competition. The MuJoCo 
physics engine~\citep{tassa2020dmcontrol} provides a high-fidelity simulation 
platform for continuous control tasks, widely adopted as a benchmark in the 
research community.

This project applies state-of-the-art DRL techniques to a 2-versus-2 football 
simulation in MuJoCo, building progressively from classical Q-Learning to the 
modern Proximal Policy Optimisation (PPO)~\citep{schulman2017proximal} 
algorithm, using the PettingZoo~\citep{terry2021pettingzoo} multi-agent 
interface and the Stable-Baselines3~\citep{raffin2021stablebaselines3} 
training library.

%%%%%%%%%%%%%%%%%%%%%%%%%%%%%%%%%%%%%%%%%%%%%%%%%%%%%%%%%%%%%%%%%%%%%%%%%%%%%%%%%%%
\subsection{Problem Statement}
\label{sec:intro_prob_art}

The core problem investigated in this project is: \textit{How can multiple 
autonomous agents learn cooperative and competitive behaviours in a physically 
realistic, continuous-action environment through reinforcement learning alone?}

Specifically, four embodied agents must learn, without any hand-crafted rules 
or human demonstrations, to:
\begin{enumerate}
    \item Navigate a continuous 3D physics space.
    \item Pursue and make contact with a moving ball.
    \item Cooperate with a teammate to advance the ball.
    \item Score goals against two opposing agents.
\end{enumerate}

This problem is difficult because rewards are sparse (only a goal scores a 
meaningful reward), the state space is high-dimensional and continuous, 
and the presence of multiple simultaneously learning agents makes the 
environment non-stationary from any individual agent's perspective.

%%%%%%%%%%%%%%%%%%%%%%%%%%%%%%%%%%%%%%%%%%%%%%%%%%%%%%%%%%%%%%%%%%%%%%%%%%%%%%%%%%%
\subsection{Aims and Objectives}
\label{sec:intro_aims_obj}

\textbf{Aim:} To design, implement, and evaluate a multi-agent deep reinforcement learning system capable of training autonomous soccer-playing agents in a physics-based 3D simulation using Proximal Policy Optimisation (PPO) with advanced reward shaping and spatial analytics.

\textbf{Objectives:}
\begin{enumerate}
    \item \textbf{O1} — Deploy and configure the MuJoCo Soccer 2v2 multi-agent environment using the Shimmy and PettingZoo frameworks, ensuring compatibility with state-of-the-art DRL architectures.
    
    \item \textbf{O2} — Design and formulate an aggressive reward shaping wrapper incorporating distance penalties, teammate coordinates, and goal-directed velocity metrics to mitigate the sparse reward challenge.
    
    \item \textbf{O3} — Train a shared cooperative team policy using PPO and the parameter-sharing paradigm to orchestrate multi-agent coordination.
    
    \item \textbf{O4} — Implement an advanced soccer analytics framework to extract ball possession statistics, teammate spatial coverage heatmaps, and coordinate-based team spreading dynamics.
\end{enumerate}

%%%%%%%%%%%%%%%%%%%%%%%%%%%%%%%%%%%%%%%%%%%%%%%%%%%%%%%%%%%%%%%%%%%%%%%%%%%%%%%%%%%
\subsection{Solution Approach}
\label{sec:intro_sol}

The solution follows a structured engineering workflow:

\begin{enumerate}
    \item \textbf{Stage 1 — Multi-Agent Environment Architecture:} We wrap the raw MuJoCo 2v2 soccer environment using PettingZoo and Shimmy compatibility layers to establish a standardized parallel interface for vector-env parallelization.
    
    \item \textbf{Stage 2 — Reward Shaping \& Core Logic:} We design a custom `SoccerRewardShapingWrapper` that translates physical states (distances, relative velocities) into immediate shaped dense rewards, guiding agents to pursue the ball and cooperate.
    
    \item \textbf{Stage 3 — Policy Optimization \& Vectorization:} We deploy a shared MultiInput PPO policy, parallelizing experience gathering over 12 vectorized environments on Google Colab with T4 GPU acceleration.
    
    \item \textbf{Stage 4 — Soccer Analytics \& Spatial Heatmaps:} We extract real-world game statistics by simulating the trained policy, computing ball possession percentages, teammate spatial coverage density heatmaps on the pitch, and team spread metrics.
\end{enumerate}

% ================================================================
\section{Project Management}
\label{sec:proj_mgmt}
% ================================================================

The project was structured over a period of approximately three weeks, following the timeline summarised in Table~\ref{tab:gantt}.

\begin{table}[!ht]
    \centering
    \caption{Project management timeline}
    \label{tab:gantt}
    \begin{tabular}{lp{5cm}c}
        \toprule
        \textbf{Phase} & \textbf{Tasks} & \textbf{Duration} \\
        \midrule
        1. Setup         & Environment installation, MuJoCo configuration, dependency management & Week 1 \\
        2. Modeling      & Reward shaping formulation, state-action analysis, PPO hyperparameters & Week 1--2 \\
        3. Main project  & Environment wrappers, training pipeline setup, Colab notebook & Week 2 \\
        4. Training      & GPU training on Google Colab, automatic checkpoint monitoring & Week 2--3 \\
        5. Analytics     & Possession analytics, teammate spacing, spatial heatmap generation & Week 3 \\
        6. Report        & LaTeX report writing and submission & Week 3 \\
        \bottomrule
    \end{tabular}
\end{table}

% ================================================================
\section{Summary}
\label{sec:intro_summary}
% ================================================================

In summary, this chapter established the context and objectives of the project.
The investigated problem --- training multiple autonomous agents to play
2-versus-2 football in a physics simulation using deep reinforcement learning ---
was described in detail. Four measurable objectives were defined, and the
three-stage solution approach was outlined. The project management plan was
also presented. Chapter~\ref{ch:lit_rev} presents a review of the existing
literature and course material relevant to this project.